% GENERATED FILE. Edit canonical lessons, then run npm run generate:books.
% canonical-source-hash: fnv1a64:f4e9ac34d01ca573
% canonical-lessons: GE-C157-tablette, GE-C157-versicherung, GE-C157-konto, GE-C157-kreditkarte, GE-C157-preis

\chapter{A tablet, Insurance, A bank account, A credit card, A price}
\label{ch:ge-a-tablet-insurance-a-bank-account-a-credit-card-a-price}
\hlchaptermodality{157}

\begin{chapteropening}
\textbf{By the end of this chapter:} \emph{I can say a tablet, insurance, a bank account, a credit card and a price.}

Complete the last lesson of chapter 157: I can say a tablet, insurance, a bank account, a credit card and a price.
\end{chapteropening}

\section[die Tablette]{die Tablette — a tablet}

\label{lesson:GE-C157-tablette}



\begin{quote}
Before the new one: say the German for a cough, then the German for the flu.
\end{quote}

\subsection*{You'll want to know: die Tablette}
\textbf{die Tablette} — \textquotedblleft{}a tablet\textquotedblright{}.

\begin{cousinweb}
\textbf{Tablette} $\leftarrow$ French \emph{tablette}, \textquotedblleft{}a little tablet\textquotedblright{} (EtymDB 2.1), a diminutive of the word that goes back to Latin \emph{tabula}, \textquotedblleft{}board\textquotedblright{} — borrowed.
\end{cousinweb}

\begin{grammarlens}[title={What you've built}]
One more word for this chapter.
\end{grammarlens}

\subsection*{Your turn}
\begin{itemize}
\raggedright
  \item \emph{Say it:} \emph{die Tablette}
  \item \emph{Say it:} \emph{die Tablette}, once more
  \item \emph{Say it:} say \emph{die Tablette}
  \item \emph{Recall:} say the German for a cough, then the German for the flu, then say \emph{die Tablette} again
\end{itemize}

\begin{tcolorbox}[breakable,colback=teal!4,colframe=teal!35!black,title={Before you move on}]
What does die Tablette mean? (\textquotedblleft{}A tablet\textquotedblright{}.) Say it once more.
\end{tcolorbox}
\section[die Versicherung]{die Versicherung — insurance}

\label{lesson:GE-C157-versicherung}



\begin{quote}
Before the new one: say the German for the flu, then the German for a tablet.
\end{quote}

\subsection*{You'll want to know: die Versicherung}
\textbf{die Versicherung} — \textquotedblleft{}insurance\textquotedblright{}.

\begin{cousinweb}
\textbf{Versicherung} is the noun of the verb \textquotedblleft{}to insure\textquotedblright{} (EtymDB 2.1), built from \emph{ver-} + \emph{sicher}, \textquotedblleft{}safe\textquotedblright{}, which goes back to Latin \emph{sēcūrus}: insurance makes you safe.
\end{cousinweb}

\begin{grammarlens}[title={What you've built}]
One more word for this chapter.
\end{grammarlens}

\subsection*{Your turn}
\begin{itemize}
\raggedright
  \item \emph{Say it:} \emph{die Versicherung}
  \item \emph{Say it:} \emph{die Versicherung}, once more
  \item \emph{Say it:} say \emph{die Versicherung}
  \item \emph{Recall:} say the German for the flu, then the German for a tablet, then say \emph{die Versicherung} again
\end{itemize}

\begin{tcolorbox}[breakable,colback=teal!4,colframe=teal!35!black,title={Before you move on}]
What does die Versicherung mean? (\textquotedblleft{}Insurance\textquotedblright{}.) Say it once more.
\end{tcolorbox}
\section[das Konto]{das Konto — a bank account}

\label{lesson:GE-C157-konto}



\begin{quote}
Before the new one: say the German for a tablet, then the German for insurance.
\end{quote}

\subsection*{You'll want to know: das Konto}
\textbf{das Konto} — \textquotedblleft{}a bank account\textquotedblright{}.

\begin{cousinweb}
\textbf{Konto} $\leftarrow$ Italian \emph{conto}, \textquotedblleft{}calculation, account\textquotedblright{} (EtymDB 2.1), from Latin \emph{computāre}, \textquotedblleft{}to reckon\textquotedblright{} — borrowed. English \emph{count} and \emph{account} share the word.
\end{cousinweb}

\begin{grammarlens}[title={What you've built}]
One more word for this chapter.
\end{grammarlens}

\subsection*{Your turn}
\begin{itemize}
\raggedright
  \item \emph{Say it:} \emph{das Konto}
  \item \emph{Say it:} \emph{das Konto}, once more
  \item \emph{Say it:} say \emph{das Konto}
  \item \emph{Recall:} say the German for a tablet, then the German for insurance, then say \emph{das Konto} again
\end{itemize}

\begin{tcolorbox}[breakable,colback=teal!4,colframe=teal!35!black,title={Before you move on}]
What does das Konto mean? (\textquotedblleft{}A bank account\textquotedblright{}.) Say it once more.
\end{tcolorbox}
\section[die Kreditkarte]{die Kreditkarte — a credit card}

\label{lesson:GE-C157-kreditkarte}



\begin{quote}
Before the new one: say the German for insurance, then the German for a bank account.
\end{quote}

\subsection*{You'll want to know: die Kreditkarte}
\textbf{die Kreditkarte} — \textquotedblleft{}a credit card\textquotedblright{}.

\begin{cousinweb}
\textbf{Kreditkarte} is \emph{Kredit}, \textquotedblleft{}credit\textquotedblright{} $\leftarrow$ French \emph{crédit} $\leftarrow$ Latin \emph{crēditum}, \textquotedblleft{}a loan, a thing trusted\textquotedblright{}, + \emph{Karte}, \textquotedblleft{}card\textquotedblright{} (EtymDB 2.1).
\end{cousinweb}

\begin{grammarlens}[title={What you've built}]
One more word for this chapter.
\end{grammarlens}

\subsection*{Your turn}
\begin{itemize}
\raggedright
  \item \emph{Say it:} \emph{die Kreditkarte}
  \item \emph{Say it:} \emph{die Kreditkarte}, once more
  \item \emph{Say it:} say \emph{die Kreditkarte}
  \item \emph{Recall:} say the German for insurance, then the German for a bank account, then say \emph{die Kreditkarte} again
\end{itemize}

\begin{tcolorbox}[breakable,colback=teal!4,colframe=teal!35!black,title={Before you move on}]
What does die Kreditkarte mean? (\textquotedblleft{}A credit card\textquotedblright{}.) Say it once more.
\end{tcolorbox}
\section[der Preis]{der Preis — a price, a prize}

\label{lesson:GE-C157-preis}



\begin{quote}
Before the new one: say the German for a bank account, then the German for a credit card.
\end{quote}

\subsection*{You'll want to know: der Preis}
\textbf{der Preis} — \textquotedblleft{}a price, a prize\textquotedblright{}.

\begin{cousinweb}
\textbf{Preis} $\leftarrow$ Middle High German \emph{prīs}, from Old French \emph{pris} $\leftarrow$ Latin \emph{pretium}, \textquotedblleft{}worth, price, reward\textquotedblright{} (EtymDB 2.1) — borrowed. English \emph{price}, \emph{prize} and \emph{praise} share it.
\end{cousinweb}

\begin{grammarlens}[title={What you've built}]
One more word for this chapter.
\end{grammarlens}

\subsection*{Your turn}
\begin{itemize}
\raggedright
  \item \emph{Say it:} \emph{der Preis}
  \item \emph{Say it:} \emph{der Preis}, once more
  \item \emph{Say it:} say \emph{der Preis}
  \item \emph{Recall:} say the German for a bank account, then the German for a credit card, then say \emph{der Preis} again
\end{itemize}

\begin{tcolorbox}[breakable,colback=teal!4,colframe=teal!35!black,title={Before you move on}]
What does der Preis mean? (\textquotedblleft{}A price, a prize\textquotedblright{}.) Say it once more.
\end{tcolorbox}
